%==============================================================
%  Burak Ömür — Resume
%  Compile with:  pdflatex burak-omur-resume.tex
%  (No exotic packages, no fontawesome, compiles anywhere.)
%==============================================================
\documentclass[a4paper,10pt]{article}

\usepackage[T1]{fontenc}
\usepackage[utf8]{inputenc}
\usepackage{charter}                 % body serif
\usepackage[scaled=0.92]{helvet}     % sans for headings
\usepackage[
  top=0.95cm, bottom=0.8cm, left=1.4cm, right=1.4cm, footskip=0.4cm
]{geometry}
\usepackage{titlesec}
\usepackage{enumitem}
\usepackage{xcolor}
\usepackage{tabularx}
\usepackage[hidelinks]{hyperref}

%---------------- palette ----------------
\definecolor{ink}{HTML}{1A1A1A}
\definecolor{accent}{HTML}{1F3A5F}
\definecolor{muted}{HTML}{5A5A5A}
\definecolor{rule}{HTML}{C9CED6}
\color{ink}

%---------------- section headings ----------------
\titleformat{\section}
  {\normalfont\sffamily\bfseries\footnotesize\color{accent}}
  {}{0em}{\MakeUppercase}[\vspace{-0.55em}{\color{rule}\rule{\linewidth}{0.5pt}}]
\titlespacing*{\section}{0pt}{0.75em}{0.45em}

%---------------- helpers ----------------
\setlist[itemize]{
  leftmargin=1.05em, itemsep=0.05em, parsep=0pt, topsep=0.2em, label=\textbullet
}
\newcommand{\sep}{\;\raisebox{0.05em}{\textcolor{rule}{\small$\vert$}}\;}

% \role{Company}{Location}{Title}{Dates}
\newcommand{\role}[4]{%
  \vspace{0.25em}%
  \noindent\begin{tabularx}{\linewidth}{@{}X r@{}}
    \textbf{#1}\,\textcolor{muted}{\small #2} & \textcolor{muted}{\small #4}\\[0.05em]
    \textit{#3} & \\
  \end{tabularx}%
  \vspace{-0.15em}%
}

% \project{Title}{Subtitle}
\newcommand{\project}[2]{%
  \vspace{0.25em}%
  \noindent\textbf{#1}\;\textcolor{muted}{\small\textit{#2}}%
  \vspace{-0.05em}%
}

% \skill{Category}{Items}
\newcommand{\skill}[2]{\noindent\textbf{\small #1:}\; {\small #2}\par\vspace{0.18em}}

\pagestyle{empty}
\raggedbottom

%==============================================================
\begin{document}

%---------------- header ----------------
\begin{center}
  {\sffamily\bfseries\Huge Burak Ömür}\\[0.32em]
  {\sffamily\large\color{accent} Senior AI Engineer}\\[0.55em]
  {\small
    Munich, DE \sep
    \href{mailto:burak.omur.1998@gmail.com}{burak.omur.1998@gmail.com} \sep
    \href{https://www.linkedin.com/in/burak-omur/}{LinkedIn} \sep
    \href{https://github.com/darktheorys}{GitHub}
  }
\end{center}
\vspace{0.4em}

%---------------- summary ----------------
\section{Profile}
\noindent AI engineer building and running production LLM systems since 2022. Lead architect of an
agentic coaching product at Retorio, owning the runtime, model routing, memory, and observability
end to end under EU data residency and GDPR constraints.

%---------------- experience ----------------
\section{Experience}

\role{Retorio GmbH}{Munich, DE}{Senior AI Engineer \textnormal{\small(Working Student until Jun 2024)}}{Mar 2022 -- Present}
\begin{itemize}
  \item Lead the architecture of \textbf{Tori}, an agentic coaching product, as the sole AI engineer
        on the team, covering runtime, model routing, state and memory, and observability.
  \item Built the runtime on a self-hosted Agent Protocol server (Aegra) with LangGraph and DeepAgents
        for long-horizon planning and subagent delegation, adding durable event logging (Redis
        Streams into PostgreSQL) and OpenTelemetry tracing into Langfuse, making multi-turn sessions
        replayable and auditable.
  \item Consolidated all model traffic behind a central LiteLLM proxy on Cloud Run against Vertex AI
        EU endpoints, adding Redis-backed fleet-wide admission control, isolating real-time and
        feedback paths, and eliminating roughly 60-second cold starts.
  \item Designed per-user memory banks and a source-addressed evidence store with bi-temporal
        provenance (PROV-O style) to support cited feedback and GDPR erasure requests.
  \item Defined the SLI and measurement approach behind a contractual SLA for an enterprise customer,
        based on real user traffic and application-layer error signals that surface mid-stream
        failures Cloud Run metrics miss.
  \item Before Tori, built FastAPI and gRPC microservices on Docker and GCP, trained in-house face
        detection and NLP models, and shipped LLM chains with retry and fallback handling.
\end{itemize}

\role{Tazi.ai}{Istanbul, TR}{Machine Learning Engineer, Intern}{Jul 2020 -- Sep 2020}
\begin{itemize}
  \item Built a feature selection method based on variational autoencoders in Scala and
        Deeplearning4j.
\end{itemize}

\role{SÜTAŞ Süt Ürünleri A.Ş.}{Istanbul, TR}{Software Engineer, Intern}{Jun 2019 -- Aug 2019}
\begin{itemize}
  \item Built an internal Android application to speed up communication between production teams
        and management.
\end{itemize}

%---------------- skills ----------------
\section{Technical Skills}
\skill{Languages}{Python, SQL, Java, C++, Scala}
\skill{LLM \& Agents}{LangGraph, LangChain, DeepAgents, Agent Protocol (Aegra), MCP / FastMCP, LiteLLM, RAG}
\skill{Observability}{Langfuse, OpenTelemetry, Grafana}
\skill{Backend \& Data}{FastAPI, gRPC, Pydantic, Docker, Redis, PostgreSQL / Cloud SQL, PgBouncer, MongoDB}
\skill{ML \& Cloud}{PyTorch, TensorFlow, JAX, computer vision; GCP (Cloud Run, Cloud SQL, Vertex AI), Azure, CI/CD}
\skill{Compliance}{GDPR, EU AI Act, EU data residency}

%---------------- projects ----------------
\section{Research}

\project{Towards Smarter and Context-Aware Retrieval Augmented Generation with LLMs}{M.Sc. Thesis, TUM}
\begin{itemize}
  \item Built a query taxonomy for RAG failure modes (feature extraction, disambiguation,
        hallucination) and compared explicit against implicit feature extraction.
  \item Created AbbrvQA, an AI-generated ambiguity dataset, and proposed a RAG pipeline with
        disambiguation and validation modules.
\end{itemize}

\project{Investigating Bias and Discrimination in Movies Using Deep Learning}{B.Sc. Thesis, Boğaziçi}
\begin{itemize}
  \item Large-scale computer vision analysis of the MovieNet dataset using pre-trained
        architectures in Python.
\end{itemize}

%---------------- education ----------------
\section{Education}
\noindent\begin{tabularx}{\linewidth}{@{}X r@{}}
  \textbf{Technical University of Munich} \textcolor{muted}{\small Munich, DE} & \textcolor{muted}{\small Oct 2021 -- Jun 2024}\\
  \textit{M.Sc. Informatics} \textcolor{muted}{\small (GPA 1.61)} & \\[0.45em]
  \textbf{Boğaziçi University} \textcolor{muted}{\small Istanbul, TR} & \textcolor{muted}{\small Sep 2016 -- Jun 2021}\\
  \textit{B.Sc. Computer Engineering} \textcolor{muted}{\small (GPA 3.75/4.00)} & \\
\end{tabularx}

\vspace{0.35em}
\noindent\textbf{\small Languages:}\; {\small Turkish (native), English (fluent), German (A2)}

\end{document}
